\documentclass[%handout %handout erstellen
]{beamer}
\usepackage[orientation=landscape,size=custom,width=16,height=9,scale=0.5,debug]{beamerposter}

\usepackage[utf8]{inputenc}
\usepackage[naustrian]{babel}
\usepackage{color}
\usepackage{url}
\usepackage{lmodern}
\usepackage{listings}
\usepackage{subfigure}

\hypersetup{pdftitle={\inserttitle},
  pdfsubject={\inserttitle},
  pdfauthor={\insertauthor},
  pdfkeywords={\inserttitle},
  pdfpagelabels={false},
%  pdfpagemode={FullScreen}
}

\definecolor{GitAddGreen}{rgb}{0, 0.5, 0} 
\definecolor{GitRemoveRed}{rgb}{0.75, 0, 0} 

\usetheme{Singapore}
%\progressbaroptions{headline=none,frametitle=picture-section,imagename=img/logo.jpg}
%\setbeamercovered{transparent}  %halbtransparente uebergaenge
%\beamertemplatenavigationsymbolsempty 	%keine nav-symbole

\mode<handout>{
  \setbeameroption{show notes}  %notizen nur bei handouts
  \usepackage{pgfpages}
  \pgfpagesuselayout{2 on 1}[a4paper, border shrink=5mm]
%  \pgfpagesuselayout{4 on 1}[a4paper, landscape, border shrink=5mm]
  \pgfpageslogicalpageoptions{1}{border code=\pgfusepath{stroke}}
  \pgfpageslogicalpageoptions{2}{border code=\pgfusepath{stroke}}
%  \pgfpageslogicalpageoptions{3}{border code=\pgfusepath{stroke}}
%  \pgfpageslogicalpageoptions{4}{border code=\pgfusepath{stroke}}
}

\title{Git Workshop}
\subtitle{VALUG - 8. April 2011}
\author[Florian Preinstorfer, Wolfgang Silbermayr]{Florian Preinstorfer, Wolfgang Silbermayr}
\date[valug - 08.04.2011]{08.04.2011}
\titlegraphic{\includegraphics[height=0.5cm]{img/git-logo-wikipedia-export}}

\begin{document}
\frame[plain]{\titlepage}

\begin{frame}{Inhalt}
  \tableofcontents[pausesections]
\end{frame}

%\include{slides/x} %section, subsection, frames
%\include{slides/y} %section, subsection, frames

\begin{frame}[plain]
  \begin{center}
    \vspace{1cm}
    Danke!
    \vspace{3cm}
    \begin{figure}[!b]
      \includegraphics[scale=0.5]{img/by-sa}
    \end{figure}
    \tiny{This work is licensed under the Creative Commons Attribution-ShareAlike 3.0 Austria license (CC-BY-SA).}
\end{center}
\end{frame}
\end{document}

% vim: tabstop=2 expandtab shiftwidth=2 softtabstop=2 autoindent
